\documentclass[letterpaper,11pt]{article}
\usepackage[top=0.8in, bottom=0.8in, left=0.9in, right=0.9in]{geometry}
\usepackage[hidelinks]{hyperref}
\usepackage{lmodern}
\usepackage{parskip}
\usepackage{microtype}
\setlength{\parskip}{6pt}
\setlength{\parindent}{0pt}
\begin{document}
\begin{flushleft}
\textbf{\large Ajay Venkatesh} \\
Brooklyn, NY \\
+1 516 585 9013 \\
\href{mailto:av3855@nyu.edu}{av3855@nyu.edu}
\end{flushleft}
\vspace{10pt}
May 21, 2026
\vspace{6pt}

Hiring Manager \\
Flex \\
New York, NY

\vspace{10pt}

Dear Hiring Manager,

My name is Ajay Venkatesh. I'm finishing my M.S. in Computer Engineering at NYU, with production software experience at PwC in Bengaluru, a summer SDE internship at Amazon, and ongoing on-campus work at NYU's Novel AI Technologies. I'm writing about the Software Engineer II, Core Platform role.

The strongest match is the backend and AWS work I owned at Amazon. I engineered services in \textbf{Java} using \textbf{Coral}, \textbf{Smithy}, \textbf{Dagger}, and \textbf{CBOR}, with low-level design patterns and unit and integration testing via \textbf{Mockito}. The infrastructure ran on \textbf{AWS CDK} provisioning \textbf{API Gateway, Lambda, DynamoDB, S3}, and \textbf{IAM}, deployed through a multi-stage and multi-region pipeline with automated promotion and rollback. Auth was wired through \textbf{OAuth, Midway SSO, and HELIS} with role-based access. That blend of disciplined Java, AWS-native infra, and careful identity work is the closest analogue I have to the Identity, Auth, Payment, and Financial services this role describes.

In parallel, at Campus Mesh I shipped a \textbf{serverless backend} on AWS using \textbf{Serverless Framework}, \textbf{Lambda (Node.js)}, \textbf{API Gateway}, \textbf{RDS/MySQL}, \textbf{DynamoDB}, \textbf{S3}, and \textbf{CloudFront} across four environments with \textbf{KMS-managed keys, SSM Parameter Store}, and \textbf{least-privilege IAM}. I built a \textbf{JWT custom authorizer} with AES-encrypted user IDs and a layered handler-service-repository-model architecture. That gives me real muscle in the same RDS, DynamoDB, ECS, Lambda, and API Gateway surface Flex runs on.

On payments specifically: at NYU I shipped a customer-facing insurance analytics platform on \textbf{React, Node.js, REST APIs} deployed via \textbf{Docker} on \textbf{EC2}, with \textbf{Stripe billing}, \textbf{RBAC}, and row-level security. I know how to reason about idempotent flows, sensitive data, and money-touching code.

On overlap with your stack: \textbf{Java}, \textbf{Python}, \textbf{MySQL}, \textbf{PostgreSQL}, \textbf{DynamoDB}, \textbf{AWS} (\textbf{ECS, Lambda, API Gateway, S3, CDK, IAM, RDS}), \textbf{Terraform}-style IaC via CDK, plus \textbf{distributed systems}, \textbf{Docker}, and \textbf{Kubernetes}. \textbf{AWS EKS} and \textbf{AWS VPN} are the two corners I haven't driven in production yet, and I'd come up the curve on them quickly inside the existing platform.

What pulls me toward Flex is the product. Rent is one of the largest, most recurring, most consequential payment categories in a renter's financial life, and reshaping the timing of those payments has direct, measurable effects on real people's monthly cash flow. Software in that category has to be correct, on time, and forgiving, and that combination of fintech rigor and real consumer outcomes is exactly the kind of work I'd like to compound my career around. The fact that it is being built here in New York is a bonus.

I'd welcome the chance to discuss further. Thanks for considering my application.

\vspace{12pt}
Sincerely, \\
Ajay Venkatesh

\end{document}
